\section{Requirements Specification}

\subsection{Game Modes}

\begin{itemize}
    \item \textbf{Two-player mode:} Two players compete on the same device.
    Black moves first, and turns alternate between Black and White after
    each valid move or pass.

    \item \textbf{Player vs.\ AI mode:} The human player controls Black,
    while the computer controls White. Before starting the game, the player
    can choose from three difficulty levels: Easy, Medium, or Hard.
    The computer responds after each human turn, with an artificial
    delay of 1500 milliseconds to simulate thinking time.
\end{itemize}

\subsection{User Interface (UI)}

\begin{itemize}
    \item A board on which Go stones can be placed.
    \item Turn indicator showing whether it's white or black's turn.
    \item Notification for stone captures. Panel for end-game scoring. A red square on board indicating Ko threat.
    \item Buttons to:
    \begin{itemize}
        \item Start a new game.
        \item Choose game mode (2-player or Player vs AI).
        \begin{itemize}
            \item Select AI difficulty level before starting a Player vs AI game.
        \end{itemize}
        \item Reset the current game to an empty board.
        \item Undo the last move.
        \item Redo the last undone move.
        \item Save the current game state to a file.
        \item Load a previously saved game state from a file.
        \item Adjust game settings, including:
        \begin{itemize}
            \item Board Color: Options to select different colors or themes for the Go board.
            \item Stone Design: Options to change the design or style of the stones.
            \item Sound Effects: Options to enable or disable sound effects for placing stones, capturing stones, end-game scoring, and background music.
            \item Background Music: Options to choose background music or disable it. Include a volume control for sound effects and music.
        \end{itemize}
        \item Exit the game.
    \end{itemize}
\end{itemize}

\subsection{Board Representation}

\begin{itemize} 
    \item Implement a graphical representation of the Go board using Raylib library.
    \item Players can place stones on the board by clicking on the desired intersection.
    \item Visual and audio feedback are provided when a stone is placed.
    \item In \textbf{2-player mode}, alternate turns between players, with the current player's turn indicated on the UI.
    \item In \textbf{Player vs AI mode}, the human player controls Black, and the computer controls White. The AI responds after each human turn, with an artificial delay of 1500 milliseconds to simulate thinking time.
\end{itemize}

\subsection{Game Logic}

\begin{itemize}
    \item Implement the rules of Go:
    \begin{itemize}
        \item Stone placement
        \item Capturing stones
        \item Ko rule
        \item Scoring: The game uses simplified Japanese-style scoring based
        on territory and captured stones, with komi added to White's score.
        It does not implement a post-game dead-stone removal phase or an
        algorithm for determining whether a group is alive or dead. 
    \end{itemize}
    \item Ensure that the game enforces valid moves and prevents illegal placements.
    \item Automatically detect and handle end-game conditions of passing two times.
\end{itemize}

\subsection{AI Opponent}

\begin{itemize}
    \item Implement three levels of AI difficulty:
    \begin{itemize}
        \item \textbf{Easy:} The AI makes quick, random moves without much evaluation.
        \item \textbf{Medium:} The AI uses the Minimax algorithm with a shallow depth, making reasonable
        moves with occasional mistakes.
        \item \textbf{Hard:} The AI uses Minimax with Alpha-Beta pruning, evaluating a deeper move tree to
        play strategically, reducing mistakes.
    \end{itemize}
    \item The AI makes moves efficiently, within a reasonable time frame, even at the Hard difficulty level.
\end{itemize}

\subsection{Game State Management}

\begin{itemize}
    \item Track the current state of the board and the game.
    \item Handle turn management, ensuring proper transitions between the player's turn and the AI's
turn or between two human players.
    \item Implement undo and redo functionality, allowing players to revert to previous game states or redo moves.
    \item Ensure that the game state is accurately updated after each move, including capturing stones and updating scores.
    \item Provide the ability to reset the game to its initial state at any time.
\end{itemize}

\subsection{Saving and Loading}

\begin{itemize}
    \item Implement functionality to save the current game state to a file in JSON format.
    \item Implement functionality to load a previously saved game state from a JSON file.
\end{itemize}

\subsection{Additional Implemented Features}

\begin{itemize}
    \item Selection of $9 \times 9$, $13 \times 13$, and $19 \times 19$ board sizes.
    \item Save and load with three slots.
\end{itemize}

\subsection{Technology Stack}

The game's core logic, AI algorithms, game state management, and graphical user interface (GUI) are implemented in C++. The application uses the Raylib library for graphics rendering, audio playback, and user input handling.

For the player-versus-computer mode, the AI uses the Minimax algorithm to determine moves at Medium and Hard difficulty levels, with Alpha-Beta pruning applied exclusively at the Hard difficulty level to optimize the decision-making process.

The game also supports saving and loading functionality using the JSON format. Game states, including stone positions, turn information, and other relevant data, are serialized into JSON files. These files can later be loaded to restore the game to its previously saved state.
